\documentclass[10pt,a5paper]{article}

% === AUTOR === (((
\author{\textit{Por Erick I. Rodríguez Juárez.}}
% )))

% === PAQUETES === (((
% \usepackage{makeidx}
% \usepackage{xltxtra}
\usepackage{amsfonts}
\usepackage{amsmath}
\usepackage{amssymb}
% \usepackage{fullpage}
\usepackage{tikz}
\usetikzlibrary{arrows.meta}
\usepackage{graphicx}
% )))

% === TIPOGRAFÍA === (((
% \setmainfont[
  % BoldFont       = bodonibi,
	% ItalicFont     = Century modern italic2.ttf,
	% BoldItalicFont = bodonibi,
	% SmallCapsFont  = lmromancaps10-regular.otf
% ]{Century_modern.ttf}
% )))

% === COMANDOS === (((
% \newcommand{\dis}{\displaystyle}
% \newcommand{\qed}{\hspace{0.5cm}\rule{0.16cm}{0.4cm}}
% \newcommand{\operator}[1]{\mathop{\vphantom{\sum}\mathchoice
% {\vcenter{\hbox{\huge $#1$}}}
% {\vcenter{\hbox{\Large $#1$}}}{#1}{#1}}\displaylimits}
% \newcommand{\suma}{\operator{\includegraphics[scale=0.09]{FOTOS/Sigma.png}}}
% \setlength{\parindent}{0mm}
% )))

% === ITALICA EN ENTORNO MATEMÁTICO === (((
% \DeclareSymbolFont{italics}{\encodingdefault}{\rmdefault}{m}{it}
% \DeclareSymbolFontAlphabet{\mathit}{italics}
% \ExplSyntaxOn
% \int_step_inline:nnnn { `A } { 1 } { `Z }
 % {  \exp_args:Nf \DeclareMathSymbol{\char_generate:nn{#1}{11}}{\mathalpha}{italics}{#1} }
% \int_step_inline:nnnn { `a } { 1 } { `z } {  \exp_args:Nf \DeclareMathSymbol{\char_generate:nn{#1}{11}}{\mathalpha}{italics}{#1}}
% \ExplSyntaxOff
% )))

\begin{document}

\titulo

\begin{enumerate}
	\item \textit{Calsifica las siguientes ecuaciones diferenciales según su tipo, orden y linealidad.} \\[2mm]
		\textbf{Solución.} 
		\begin{table}[ht]
			\centering
			\begin{tabular}{|*{6} {l|}}
				\hline
				Ec. & Tipo & Orden & Lineal & V.S. \\ \hline
				\(\sqrt{1-x} \cdot \dfrac{d^2y}{dx^2} +2x \cdot \dfrac{dy}{dx} = y\) & Ordinaria & \(2^{do}\) Orden & Lineal & No \\ \hline
				\(y'' -t^2yy' = \sin t\) & Ordinaria & \(2^{do}\) Orden & No lineal & No. \\ \hline
			\end{tabular}
		\end{table}
	\item \textit{Verifica que la función \(y(x) = 3 \sin 2x+e^{-x}\) es solución explícita de la ecuación diferencial} 
		\[
			y'' +4y=5e^{-x}.
		\]
		\textbf{Solución.} Tenemos que
		\[
			\begin{array}{rcl}
				y' & = & 6 \cos 2x-e^{-x} \\[2mm]
				y'' & = & -12 \sin 2x+e^{-x}
			\end{array}
		\]
		Así que
		\[
			(-12 \sin 2x+e^{-x}) +4(3 \sin 2x+e^{-x}) = 5e^{-x}.
		\]
	\item \textit{Compruebe que \(y- \ln y=x^2+1\) es solución general de la ecuación diferencial}
		\[
			y' = \dfrac{2xy}{y-1}.
		\]
		\textbf{Solución.} Tenemos que
		\[
			\begin{array}{rcl}
				y' +2 \dfrac{y'}{y} =2x & \;\implies\; & y' \Bigg(1+ \dfrac{2}{y}\Bigg) = 2x \\[2mm]
				& \;\implies\; & y' = \dfrac{2x}{1+2/y} = \dfrac{2xy}{y+2}. \qed 
			\end{array}
		\]
	\item \textit{Determina si las siguientes ecuciones diferneciales son de V.S. o lineales y resuélvelas con el método correspondiente. En cada inciso, indica cuál es solución general implícita y encuentra la explícita.} 
		\begin{enumerate}
			\item \(x \cdot \dfrac{dy}{dx} -2y=x^3e^{-x} -5x\). \\[2mm]
				\textbf{Solución.} \fbox{Es lineal. No es de V.S.} \\[2mm]
				Tenemos que
				\[
					\dfrac{dy}{dx} - \dfrac{2}{x} y = x^2e^{-x} -5.
				\]
				Entocnes \(p(x) = - \dfrac{2}{x}\), luego
				\[
					\mu (x) = e^{\int p} = e^{\int -2/x} = e^{-2  \log x} =e^{\log x^{-2}} =x^{-2}.
				\]
				Por tanto,
				\[
					\begin{array}{rcl}
						x^{-2} \dfrac{dy}{dx} - \dfrac{2}{x^3} y & = & e^{-x} -5x^{-2} \\[5mm]
						\dfrac{d}{dx} \Bigg(x^{-2} y\Bigg) & = & e^{-x} -5x^{-2} \\[5mm]
						\dis\int \dfrac{d}{dx} \Bigg(x^{-2} y\Bigg) dx & = & \dis\int (e^{-x} -5x^{-2} )dx \\[5mm]
						yx^{-2} & = & -e^{-x} + \dfrac{5}{x} +C, \hspace{1cm} C \in \mathbf{R} \\[5mm]
						y & = & -x^2e^{-x} + 5x+ Cx^2, \hspace{1cm} C \in \mathbf{R}.
					\end{array}
				\]
			\item \((x+xy^2) dx+e^{x^2} ydy=0\). \\[2mm]
				\textbf{Solución.} \fbox{La ecuación es de V.S.} \\[2mm]
				Tenemos que
				\[
					\begin{array}{rcl}
						x(1+y^2) dx +e^{x^2} ydy & = & 0 \\[5mm]
						e^{x^2} ydy & = & -x(1+y^2) dx \\[5mm]
						\dfrac{ydy}{1+y^2} & = & \dfrac{-xdx}{e^{x^2}} \\[5mm]
						\dis\int \dfrac{ydy}{1+y^2} & = & \dis\int \dfrac{-xdx}{e^{x^2}} \\[5mm]
						\dfrac{1}{2} \dis\int \dfrac{du}{1+u} & = & -\dfrac{1}{2} \dis\int \dfrac{dw}{e^{w}} \\[5mm]
					\end{array}
				\]
				\[
					\begin{array}{rcl}
						\log (1+u) & = & +e^{-w} +C \\[5mm]
						\log (1+y^2) & = & e^{-x^2} +C \\[5mm]
						1+y^2 & = & e^{e^{-x^2} +C} \\[5mm]
						y & = & \pm \sqrt{e^{e^{-x^2} +C} -1} \;,\; \hspace{1cm} C \in \mathbf{R}.
					\end{array}
				\]
		\end{enumerate}
	\item \textit{Resuelve la siguiente ecuación diferencial \(x \cdot \dfrac{dy}{dx} = xy \cos x-2xy\), y determina la solución particular explícita que pasa por el punto \((\pi ,1)\).} \\[2mm]
		\textbf{Solución.} Dividamos entre \(x\), todos los términos.
		\[
			\begin{array}{rcl}
				\dfrac{dy}{dx} &=& y \cos x-2y \\[2mm]
				\dfrac{dy}{dx} -y(\cos x-2) & = & 0.
			\end{array}
		\]
		Entonces, se tiene \(p(x) = -\cos x+2\). Luego
		\[
			\mu (x) = e^{\int p} = e^{\int (-\cos x+2) dx} = e^{-\sin x+2x}.
		\]
		Por tanto,
		\[
			\begin{array}{rcl}
				e^{\sin x-2x} \dfrac{dy}{dx} -ye^{- \sin x+2x} (\cos x-2) & = & 0 \\[5mm]
				\dfrac{d}{dy} \big(ye^{- \sin x+2x}\big) & = & 0\\[5mm]
				ye^{- \sin x+2x} & = & k, \hspace{1cm} k \in \mathbf{R} \\[5mm]
				y & = & ke^{\sin x-2x} , \hspace{1cm} k \in \mathbf{R}.
			\end{array}
		\]
		Ahora, se requiere que \(y(\pi) =1\), entonces
		\[
			1 = y(\pi) = ke^{\sin \pi -2 \pi} = ke^{-2 \pi}.
		\]
		Entonces, \(k=e^{2 \pi}\). La solución del P.V.I. \(y(\pi) =1\), es
		\[
			y = e^{2 \pi} \cdot e^{\sin x-2x}.
		\]
\end{enumerate}

\end{document}
